\subsection{Graph product polynomial optimization}
\label{sec:gpop}
Consider three groups($A,B,AB$) where each group has two variables and so, $X=\{X_{j,i} : i \in \{A,B,AB\},j \in \{1,2\}\}$. All the variables of group $A$ commute with all the variables of group $B$ but not with the variables of group $AB$. The variables $X_1^i,X_2^i$ of each group $i$ follow the constraints as in problem~\ref{ncpop_pna}. The optimization problem we consider reads as follows.

\begin{equation}
  \label{gppop}
\begin{split}
       p^* &=\underset{\underline{X} , \psi , \mathcal{H}}{inf} \: \braket{\psi| \:\bigl[\sum_{i \in \{A,B,AB\}} X_{1,i} \mult X_{2,i} + X_{2,i} \mult X_{1,i} + \sum_{i,j \in \{A,B,AB\}} X_{1,i}*X_{2,j} \:\bigr]\: |\psi}, \\
    & s.t \quad \qquad X_{1,i}^2 - X_{1,i} = 0,\forall i \in \{A,B,AB\} \\
    & \quad \qquad [X_{i,A},X_{j,B}]=0, \forall i,j \in \{1,2\}\\
    & \qquad \quad -X_{2,i}^2+X_{2,i}+1/2 \geq 0,\forall i \in \{A,B,AB\} \\
    & \qquad \quad \braket{\psi , \psi} = 1. \\ 
\end{split}
\end{equation} 
With PCPOP, the above problem can be solved as follows.
\begin{code}
    # Monoid fro group A
    @ncmonoid A X1A X2A
    # Monoid for group B
    @ncmonoid B X1B X2B
    # Monoid for group AB
    @ncmonoid AB X1AB X2AB
    vars=[X1A,X2A,X1B,X2B,X1AB,X2AB]
     #All variables are hermitian and projectors. This constraint only involve variables of same group.
    Projector.(vars)
    # We build each child monoids.
    # We create the parent monoid using the child monoids.
    M=graph_product_monoid("M",[A,B,AB])
    # We enforce the commutation relations between the child monoids, this is enough to enforce all the required commutation relations.
    @comms A B
    # Build the parent monoid, it automatically builds the child monoids.
    build(M)
    # Define the objective function.
    obj_local=X1A*X2A + X2A*X1A + X1B*X2B + X2B*X1B + X1AB*X2AB + X2AB*X1AB
    obj_cross=X1A*X2B + X1A*X2AB + X1B*X2A + X1B*X2AB + X1AB*X2A + X1AB*X2B
    obj= obj_local + obj_cross
    # Define operator inequalities.
    op_ge=[i-i^2+0.5 for i in vars]
    # Define  level of relaxation.
    level=2
    # Solve the graph product polynomial optimization problem.
    ov,model,var_dict,principal_matrix=npa(obj,level;op_ge=op_ge)
    println("Optimal value is ",ov)
\end{code}

The steps to solve a graph product polynomial optimization problem are as follows.
\begin{itemize}
    \item First we define the child monoids. One can use already-defined macros like \texttt{@ncmonoid} , \texttt{@pcmonoid} or create custom monoids. 
    \item Then we enforce operator equality constraints on the child monoids and build them.
    \item Then we create the parent monoid using the built-in function \texttt{graph\_product\_monoid} which takes as input, the name of the parent monoid and an array of child monoids.
    \item Then we define the commutation relations between the child monoids using the built-in macro \texttt{@comms}.
    \item Then we build the parent monoid using \texttt{build(M)}.
    \item Finally, we call the function \texttt{npa} to solve the optimization problem. You can give operator equalities that act on variables from different child monoids using the optional argument $\texttt{op\_eq}$ of \texttt{npa} function. As before, you can also give operator inequalities(\texttt{op\_ge}), expectation value equalities(\texttt{tr\_eq}) and inequalities(\texttt{tr\_ge}) using the respective optional arguments of \texttt{npa} function.
\end{itemize}
One can simplify the above code using the built-in macro \texttt{@subsystem} as follows.
\begin{code}
    # Directly create parent monoid M with child monoids A,B,AB each having two hermitian variables.
    @subsystem M A[2,0] B[2,0] AB[2,0]
    # All the variables of parent monoid can be accessed using variables(M). The variables are {a1,a2,b1,b2,ab1,ab2}.
    vars=variables(M)
    # Enforcing projector constraint on all the variables. These constraints are handled using Grobner basis reduction. 
    Projector.(vars)
    # Building parent parent monoid, automatically build child monoids.
    build(M)
    # Enforcing operator inequalities on all the variables.
    op_ge=[i-i^2+0.5 for i in vars]
    # Define objective function.
    obj_local = a[1]*a[2] + a[2]*a[1] + b[1]*b[2] + b[2]*b[1] + ab[1]*ab[2] + ab[2]*ab[1]
    obj_cross = a[1]*b[2] + a[1]*ab[2] + b[1]*a[2] + b[1]*ab[2] + ab[1]*a[2] + ab[1]*b[2]
    obj = obj_local + obj_cross
    # Define level of relaxation.
    level=2
    # Solve the graph product polynomial optimization problem.
    ov,model,var_dict,principal_matrix=npa(obj,level;op_ge=op_ge)
    println("Optimal value is ",ov)
\end{code}